\documentclass{article}
\usepackage{amsmath, amssymb, amsthm}
\usepackage{hyperref}
\usepackage{geometry}
\geometry{margin=1in}

\title{Erd\H{o}s Problem \#348}
\date{Last updated: 2025-08-31}
\author{Source: erdosproblems.com}

\begin{document}
\maketitle

\noindent\textbf{Status:} OPEN \quad \textbf{No prize}

\noindent\textbf{Tags:} number theory, complete sequences

\medskip

\noindent\textbf{Problem Statement:}

For what values of $0\leq m&#60;n$ is there a complete sequence $A=\{a_1\leq a_2\leq \cdots\}$ of integers such that {UL} {LI} $A$ remains complete after removing any $m$ elements, but {/LI} {LI} $A$ is not complete after removing any $n$ elements? {/LI} {/UL}




The Fibonacci sequence $1,1,2,3,5,\ldots$ shows that $m=1$ and $n=2$ is possible. The sequence of powers of $2$ shows that $m=0$ and $n=1$ is possible. The case $m=2$ and $n=3$ is not known. van Doorn has shown that no such sequence exists for $2\leq m&lt;n$ if we interpret complete in the strong sense that\[\left\{ \sum_{n\in B}n : \textrm{ for all finite }B\subset A\right\}=\mathbb{N}.\]Erd\H{o}s and Graham most likely meant the weaker notion of completeness which allows finitely many exceptions, however.


\bigskip
\noindent\small{Source: \url{https://www.erdosproblems.com/348}}
\end{document}
